% Simon Röhrl — Curriculum Vitae
% Classic black-and-white two-page CV, article class, Computer Modern.
\documentclass[10pt,a4paper]{article}

\usepackage[utf8]{inputenc}
\usepackage[T1]{fontenc}
\usepackage[a4paper,margin=1.8cm]{geometry}
\usepackage{enumitem}
\usepackage[colorlinks=true,urlcolor=blue,linkcolor=black,citecolor=black]{hyperref}
\urlstyle{same}

\pagestyle{empty}
\setlength{\parindent}{0pt}
\emergencystretch=1.5em

% Section heading: small caps headline with a thin rule underneath.
\newcommand{\cvsection}[1]{%
  \vspace{5pt}%
  {\large\scshape #1}\par
  \vspace{1pt}\hrule\vspace{4pt}%
}

% Entry: date column left, description right (classic German CV convention).
\newlength{\datewidth}
\setlength{\datewidth}{2.9cm}
\newcommand{\cventry}[2]{%
  \begin{minipage}[t]{\datewidth}\raggedright #1\end{minipage}%
  \hfill
  \begin{minipage}[t]{\dimexpr\textwidth-\datewidth-0.3cm\relax}#2\end{minipage}%
  \par\vspace{4pt}%
}

\setlist[itemize]{leftmargin=1.1em,itemsep=0.5pt,parsep=0pt,topsep=1pt,label=--}

\begin{document}

%----------------------------------------------------------------------
% Header
%----------------------------------------------------------------------
\begin{center}
  {\LARGE\bfseries Simon Röhrl}\\[3pt]
  Am Herrnberg 5, 93158 Teublitz, Germany
  \;\textbar\; \href{mailto:simon.roehrl01@gmail.com}{simon.roehrl01@gmail.com}
  \;\textbar\; \href{https://simonr.net}{simonr.net}\\[2pt]
  Google Scholar: \href{https://scholar.google.com/citations?user=UL63m84AAAAJ}{scholar.google.com/citations?user=UL63m84AAAAJ}\\[2pt]
  ORCID: \href{https://orcid.org/0009-0006-0768-0485}{orcid.org/0009-0006-0768-0485}
  \;\textbar\; GitHub: \href{https://github.com/simonroehrl}{github.com/simonroehrl}
\end{center}

\vspace{2pt}
Founder \& full-stack engineer and PhD researcher working on eye tracking, web technologies and AI. Founder of RSites GmbH; doctoral researcher at the University of Eastern Finland with 20+ peer-reviewed publications. Builds end-to-end, works daily with Claude Code and Codex, and takes ideas to working products in days.

%----------------------------------------------------------------------
\cvsection{Education}
%----------------------------------------------------------------------

\cventry{03/2025 -- present}{%
  \textbf{PhD (Doctoral Researcher)}, University of Eastern Finland, Joensuu\\
  Research on eye tracking and the web.}

\cventry{10/2023 -- 03/2025}{%
  \textbf{Master of Applied Research in Engineering Sciences (M.Sc.)}, OTH Regensburg --- final GPA 1.0 (every graded module ``sehr gut'', 1.0)\\
  Selected coursework (all 1.0): Deep Learning; Functional Safety \& IT Security; Probability, Statistics \& Optimization; Digitalization \& Ethics; Eye Tracking in Engineering Sciences; Agile Software Development with Scrum; research methods (grant writing, scientific presenting, start-up case study).}

\cventry{10/2019 -- 09/2023}{%
  \textbf{B.Eng.\ Electrical Engineering \& Information Technology}, OTH Regensburg\\
  Dual study programme with OSRAM.}

\cventry{09/2011 -- 06/2019}{%
  \textbf{Abitur}, Johann-Michael-Fischer-Gymnasium Burglengenfeld.}

%----------------------------------------------------------------------
\cvsection{Professional Experience}
%----------------------------------------------------------------------

\cventry{08/2023 -- present}{%
  \textbf{Founder}, RSites GmbH, Teublitz (\href{https://rsites.net}{rsites.net})
  \begin{itemize}
    \item Corporate websites built accessibility-first (European Accessibility Act ready); 30+ websites delivered overall.
    \item Recent frontends with agency partner Gessulat + Gessulat: Suzuki (corporate \& careers), Lürssen, LKQ Europe, Teledata redesign, Wurzer/Wumas, Servero, VDW; other brands incl.\ IG Metall (BMW, Siemens).
    \item Design-to-code tooling and rendering infrastructure.
    \item Responsible for delivery, finance, accounting, legal/tax, marketing and hiring; manages a small freelancer team (freelancers employed for 3+ years).
  \end{itemize}}

\cventry{03/2025 -- present}{%
  \textbf{Research Associate, project REACH}, OTH Regensburg
  \begin{itemize}
    \item Eye-tracking glasses and web usability with industry partners (WaveEye, Optomech).
    \item Eye-tracker control via WebUSB; accuracy study with 34 participants using a frame-accurate evaluation method.
    \item Usability study (25 participants, within-subjects) of a new eye-tracking-glasses prototype.
  \end{itemize}}

\cventry{10/2023 -- 03/2025}{%
  \textbf{Research Associate, project HASKI}, OTH Regensburg (haski-learning.de)
  \begin{itemize}
    \item Full-stack development of an eye-tracking web platform (Node.js backend, Docker/Kubernetes, GitLab CI/CD); designed and ran 5+ eye-tracking usability studies.
    \item Lab IT administration (Linux servers, Bash, Synology, DNS, certificates); supervised 10+ Bachelor's/Master's/project theses.
    \item Co-led funding applications and international cooperation agreements leading to funded projects, incl.\ institutional cooperation Abertay University (Scotland) -- OTH Regensburg.
    \item Open-source Moodle plugin for AI learning paths (PHP/JS, SQL); session recording with rrweb for learning analytics; first-author paper on ethics of adaptive AI learning systems.
  \end{itemize}}

\cventry{03/2022 -- 09/2023}{%
  \textbf{Junior Web Developer}, Gessulat + Gessulat GmbH \& Co.\ KG
  \begin{itemize}
    \item Delivered customer-facing sites end-to-end.
    \item Led pixel-perfect frontend implementation for Aldi Süd IT Jobs, Mytheresa Careers, Akkodis, LKQ Europe Careers, Amapharm, Simplicity and Globl.Contact.
  \end{itemize}}

\cventry{05/2021 -- 02/2022}{%
  \textbf{Freelance Web Developer} (self-employed)
  \begin{itemize}
    \item Built and owned the full Shopify stack for Bears with Benefits (beauty/nutrition start-up, products at dm; later exited).
    \item Shops, server-side tagging (GA + Facebook Conversions API), GDPR consent management, performance/SEO, landing pages, technical A/B testing.
  \end{itemize}}

\cventry{07/2019 -- 01/2022}{%
  \textbf{Dual Study \& Apprenticeship (Elektroniker für Betriebstechnik)}, OSRAM Opto Semiconductors GmbH, Regensburg --- completed apprenticeship 01/2022, alongside the dual B.Eng.\ programme.}

%----------------------------------------------------------------------
\cvsection{Research \& Academic Activities}
%----------------------------------------------------------------------

\begin{itemize}
  \item Chair and main organiser of the METR workshop (Metrics in Eye Tracking Research) at ETRA 2026, Marrakech, Morocco --- the leading eye-tracking conference: 7 workshop papers, own peer-review process, programme committee lead; community of $\sim$68 researchers forming a working group on metric standardisation.
  \item ETRA 2026 paper ``Computational Operationalization Choices for Spatial Data Quality in Screen-Based Eye Tracking'': 8 usually-implicit calculation decisions for accuracy and precision, e.g.\ exact geometric median vs coordinate-wise approximation: $0.198^{\circ}$ vs $0.146^{\circ}$ ($\sim$35\% difference on identical data).
  \item 20+ peer-reviewed papers (first and co-author) on eye tracking, web technologies and AI; mathematical-theoretical and literature work on eye-tracking metrics on the web.
  \item Research stays in Finland, Morocco, the Netherlands and Spain; international collaborators at Utrecht, UEF Joensuu, Lund, Harvard, Clemson, Nebraska, Abertay (Scotland) and Universiti Malaya; Uni Stuttgart and Uni Saarland nationally.
  \item Supervised 10+ Bachelor's/Master's/project theses; co-led funding applications.
\end{itemize}

%----------------------------------------------------------------------
\cvsection{Selected Publications}
%----------------------------------------------------------------------

\begin{enumerate}[leftmargin=1.6em,itemsep=1pt,parsep=0pt,topsep=2pt]
  \item S.~Röhrl, R.~Bednarik, J.~H.~Mottok. \emph{Computational Operationalization Choices for Spatial Data Quality in Screen-Based Eye Tracking}. Proceedings of the 2026 Symposium on Eye Tracking Research and Applications (ETRA), 2026.
  \item S.~Röhrl, S.~Staufer, F.~Bugert, V.~K.~Nadimpalli, F.~Hauser, L.~Grabinger. \emph{Ethical Considerations of AI in Education: A Case Study Based on Pythia Learning Enhancement System}. IEEE Access, 2025.
  \item S.~Röhrl, F.~Hauser, T.~Ezer, L.~Grabinger, J.~Mottok. \emph{WebGazeTrack: A Web-Based Eye Tracking Tool}. Proceedings of the 6th European Conference on Software Engineering Education, 2025.
  \item S.~Röhrl, S.~Staufer, V.~K.~Nadimpalli, F.~Bugert, F.~Hauser, L.~Grabinger. \emph{Pythia-AI Suggested Individual Learning Paths for Every Student}. INTED2024 Proceedings, 2024.
  \item T.~Ezer, F.~Hauser, S.~Röhrl, J.~H.~Mottok. \emph{Inducing Global and Focal View in UML Class Diagram Comprehension}. Proceedings of the 2026 Symposium on Eye Tracking Research and Applications (ETRA), 2026.
\end{enumerate}
\vspace{2pt}
Full list: Google Scholar profile (\href{https://scholar.google.com/citations?user=UL63m84AAAAJ}{scholar.google.com/citations?user=UL63m84AAAAJ}).

%----------------------------------------------------------------------
\cvsection{Projects}
%----------------------------------------------------------------------

\begin{itemize}
  \item \textbf{Blitzframes} (currently building) --- accelerated video rendering for Remotion: browser rasterisation, JS execution and H.264 encoding on dedicated hardware, aimed at teams with real rendering costs.
  \item \textbf{WebGazeTrack} --- Chrome extension driving Tobii eye-tracking hardware directly via WebUSB (Tobii Pro SDK), plug-and-play, dynamic AOIs from DOM/CSS selectors (ECSEE 2025 paper).
  \item \textbf{LaS3 online study software} --- web-based eye-tracking study platform (create/record/analyze, slideshow designer, webcam gaze recording, re-implemented Tobii I-VT fixation filter).
  \item \textbf{Dynamic AOI tracking with LLMs} --- hybrid DOM structure + screenshots + LLM annotation to track moving web content pixel- and frame-accurately (ETRA 2026, ``Tracing Dynamic AOIs on the Web'').
  \item \textbf{frontend-builder} (2026) --- automated Figma-to-frontend pipeline generating the implemented frontend from the design file; built for a Munich client, in production.
  \item \textbf{ecb-rates.com} --- derives expected ECB rate paths from futures pricing (work in progress).
  \item \textbf{Moodle plugin for AI learning paths} --- open source (PHP/JS).
\end{itemize}

%----------------------------------------------------------------------
\cvsection{Technical Skills}
%----------------------------------------------------------------------

\begin{itemize}
  \item \textbf{Web:} TypeScript, React, Node.js, Remotion, rrweb, D3, Chart.js, WordPress/PHP, TYPO3, Shopify
  \item \textbf{Data/ML:} Python, NumPy, pandas, SciPy, matplotlib, TensorFlow
  \item \textbf{Infrastructure:} Docker, Kubernetes, GitLab CI/CD, Linux, Bash, Synology, DNS/TLS certificates
  \item \textbf{Agents/AI tooling:} Claude Code, Codex, MCP, Playwright/CDP
  \item \textbf{Analytics:} server-side tagging, GA \& Facebook Conversions API, GDPR consent, SEO, performance
  \item \textbf{Methods:} Git, Scrum, agile development
\end{itemize}

%----------------------------------------------------------------------
\cvsection{Awards \& Activities}
%----------------------------------------------------------------------

\begin{itemize}
  \item Mathematics: 5th place Bavaria / 27th nationally, Pangea competition (10{,}000+ entrants); 1st prize Landeswettbewerb Mathematik 2014 and 2015; 1st prize Känguru 2013 and 2015; Fürther Mathematik-Olympiade (very good success, winners' seminar at OTH Regensburg).
  \item 1st Landespreis, Bundeswettbewerb Latein; winner Ostbayerisches ScienceCamp 2017.
  \item Athletics: current Bavarian runner-up, 400\,m hurdles, 2026; 2x Bavarian youth runner-up (400\,m, pentathlon) 2018/19; 2x North Bavarian champion (pentathlon) 2018/19.
  \item Languages: German (native), English (fluent).
\end{itemize}

\end{document}
